\documentclass[aspectratio=169]{beamer}
\usetheme{Madrid}
\usecolortheme{whale}

\usepackage[utf8]{inputenc}
\usepackage[spanish]{babel}
\usepackage{amsmath,amssymb}
\usepackage{graphicx}
\usepackage{listings}
\usepackage{booktabs}

\lstset{
    language=Python,
    basicstyle=\ttfamily\footnotesize,
    keywordstyle=\color{blue}\bfseries,
    stringstyle=\color{red},
    commentstyle=\color{gray}\itshape,
    breaklines=true,
    showstringspaces=false
}

\title[INT210 - Sesión 0.3]{Sesión 0.3: Bucles Iterativos Deterministas (\texttt{for}) y Comprensiones}
\subtitle{Módulo 0: Fundamentos Algorítmicos en Python}
\author[Diplomado en Ciencia de Datos]{Cátedra de Inteligencia Artificial y Ciencia de Datos}
\institute[INT210]{Machine Learning From Scratch}
\date{Semana 1}

\begin{document}

\frame{\titlepage}

\begin{frame}{Contenido de la Sesión}
    \tableofcontents
\end{frame}

\section{La Cinta Transportadora e Iteradores}

\begin{frame}{La Analogía Infantil: La Cinta Transportadora}
    \begin{columns}
        \column{0.55\textwidth}
        \begin{itemize}
            \item El robot no acumula todas las galletas en su mesa; recibe \textbf{una a la vez}.
            \item Procesa el elemento actual y solicita el siguiente mediante una señal.
            \item Al terminarse la bandeja, el flujo se detiene de forma limpia.
        \end{itemize}
        \column{0.45\textwidth}
        \begin{alertblock}{Evaluación Perezosa (\textit{Lazy Evaluation})}
            Calcular elementos bajo demanda ahorra memoria RAM y permite procesar volúmenes masivos de datos (*Big Data*).
        \end{alertblock}
    \end{columns}
\end{frame}

\section{Protocolo Iterador y Memoria RAM}

\begin{frame}{El Protocolo Iterador de Python}
    \begin{center}
        \texttt{coleccion.\_\_iter\_\_()} $\longrightarrow$ \texttt{iterador.\_\_next\_\_()} $\longrightarrow$ \texttt{StopIteration}
    \end{center}
    \vspace{0.3cm}
    \begin{columns}
        \column{0.5\textwidth}
        \textbf{Secuencia Virtual (\texttt{range}):}
        \begin{itemize}
            \item Almacena solo \texttt{(start, stop, step)}.
            \item \texttt{sys.getsizeof(range(10\^{}8))} $\approx 48$ bytes.
            \item Complejidad espacial: $\mathcal{O}(1)$.
        \end{itemize}
        \column{0.5\textwidth}
        \textbf{Lista Materializada en RAM:}
        \begin{itemize}
            \item Reserva memoria contigua para cada puntero.
            \item \texttt{list(range(10\^{}8))} $\approx 800$ MB - 3.5 GB.
            \item Complejidad espacial: $\mathcal{O}(N)$.
        \end{itemize}
    \end{columns}
\end{frame}

\section{Comprensiones Pythonicas y Enumerate}

\begin{frame}[fragile]{List Comprehensions y Transformación Matricial}
    \textbf{Transformación de Píxeles RGB a Escala de Grises:}
\begin{lstlisting}
imagen_grises = [
    [int(0.299*r + 0.587*g + 0.114*b) for r, g, b in fila]
    for fila in imagen_rgb
]
\end{lstlisting}
    \vspace{0.2cm}
    \textbf{Uso de \texttt{enumerate} para indexación bidimensional:}
\begin{lstlisting}
for y, fila in enumerate(imagen_rgb):
    for x, pixel in enumerate(fila):
        procesar_coordenada(x, y, pixel)
\end{lstlisting}
\end{frame}

\section{Casos de Uso en Ciberseguridad y Arte}

\begin{frame}[fragile]{Filtrado de Puertos de Red con Comprensiones}
\begin{lstlisting}
PUERTOS_CRITICOS = {23, 445, 3389}

alertas = [
    log for log in logs_escaneo 
    if log['puerto'] in PUERTOS_CRITICOS and log['estado'] == 'Abierto'
]
\end{lstlisting}
\begin{exampleblock}{Ventajas del Paradigma Declarativo}
    Mayor expresividad sintáctica, ejecución interna optimizada en C por el intérprete y código libre de efectos secundarios (*side-effects*).
\end{exampleblock}
\end{frame}

\begin{frame}{Conclusiones}
    \begin{itemize}
        \item El bucle \texttt{for} en Python es un consumidor de iteradores de alto nivel.
        \item Las listas por comprensión permiten vectorizar la lógica de transformación de matrices y registros.
        \item \textbf{Siguiente Sesión (0.4):} Bucles Indeterminados (\texttt{while}), Señales Centinela y Convergencia Numérica.
    \end{itemize}
\end{frame}

\end{document}
